\documentclass[conference]{IEEEtran}

\usepackage{graphicx}
\usepackage{amsmath}
\usepackage{booktabs}
\usepackage{cite}
\usepackage{url}

\title{Flight Delay Prediction Using Logistic Regression and Gradient Boosting}

\author{
\IEEEauthorblockN{Atharva Tiwari}
\IEEEauthorblockA{
Flight Delay Prediction Project\\
India
}
}

\begin{document}

\maketitle

\begin{abstract}
Flight delays affect passengers, airlines, and airport operations.
This project develops a machine learning pipeline to predict flight
arrival delays using historical flight data. The study uses the
2013 New York City flights dataset and defines a delayed flight as
one with an arrival delay of at least 15 minutes. Logistic Regression
is used as a baseline model, while XGBoost is used as the gradient
boosting model. The workflow includes exploratory data analysis,
feature engineering, chronological data splitting, model evaluation,
and performance comparison. The results show that XGBoost achieves
a higher precision-recall area under the curve and recall than the
baseline at the evaluated thresholds, although both models have
limited discrimination performance.
\end{abstract}

\begin{IEEEkeywords}
Flight delay prediction, machine learning, Logistic Regression,
XGBoost, feature engineering, classification
\end{IEEEkeywords}

\section{Introduction}

Flight delays create inconvenience for passengers and operational
challenges for airlines. Predicting delays in advance can help
stakeholders understand the likelihood of disruption and support
better operational planning.

Machine learning methods can identify patterns in historical flight
data and use them to estimate the probability of a future delay.
This project implements a classification pipeline using historical
flight records, engineered temporal features, and two machine learning
models.

The current work focuses on the first phase of the project:
exploratory analysis, feature engineering, model training, and
evaluation.

\section{Problem Formulation}

The task is formulated as a binary classification problem. Let
$A_i$ represent the arrival delay, in minutes, for flight $i$.
The target variable is defined as

\begin{equation}
y_i =
\begin{cases}
1, & A_i \geq 15,\\
0, & A_i < 15.
\end{cases}
\end{equation}

Records with missing arrival-delay labels are removed before
constructing the target. The model uses information about the
scheduled flight and route to predict whether the flight will be
delayed.

Actual arrival delay and other post-departure information are
excluded from the input features to reduce target leakage.

\section{Dataset and Preprocessing}

The project uses the 2013 New York City flights dataset, which
contains flight records associated with New York-area airports.

The raw dataset contains 336,776 rows. After removing records with
missing arrival-delay labels, 327,346 rows remain for analysis.
The resulting processed dataset contains 18 columns: one scheduled
datetime column, 16 engineered input features, and the binary target.

The observed positive-class rate is approximately 24.47\%.

Preprocessing includes removing records without target labels,
constructing the binary delay target, converting scheduled time
information into datetime values, and arranging observations
chronologically. A chronological split is used, with 70\% of the
data for training, 15\% for validation, and the final 15\% for testing.

\section{Feature Engineering}

Features are designed to represent scheduled timing and flight
characteristics. They include month, day, hour, minute, day of week,
day of year, and flight distance. The carrier, origin airport, and
destination airport are also included.

Cyclical time features represent periodic relationships. For example,
the hour of day is encoded as

\begin{equation}
h_{\sin} = \sin\left(\frac{2\pi h}{24}\right),
\end{equation}

\begin{equation}
h_{\cos} = \cos\left(\frac{2\pi h}{24}\right).
\end{equation}

Similar sine and cosine transformations are used for day of week
and day of year. These representations allow values near the end
and beginning of a cycle to remain close in feature space.

\section{Methodology}

\subsection{Logistic Regression}

Logistic Regression is used as the baseline classifier. It estimates
the probability of the positive class using a logistic function:

\begin{equation}
P(y=1 \mid \mathbf{x}) =
\frac{1}{1+\exp(-(\mathbf{w}^{T}\mathbf{x}+b))}.
\end{equation}

A classification threshold converts the estimated probability into
a binary prediction.

\subsection{XGBoost}

XGBoost is a gradient boosting method that combines multiple decision
trees. Each successive tree contributes to improving the model's
objective. It can capture nonlinear relationships and interactions
between flight timing, distance, carrier, and airport features.

The project uses 400 estimators, a maximum tree depth of 6, a learning
rate of 0.05, subsampling of 0.8, and column subsampling of 0.8.
The binary logistic objective is used for classification.

\subsection{Evaluation Metrics}

Model performance is assessed using precision, recall, F1-score,
accuracy, area under the precision-recall curve (PR-AUC), area under
the receiver operating characteristic curve (ROC-AUC), and the
Brier score.

Precision measures the fraction of predicted delays that are actual
delays. Recall measures the fraction of actual delays identified by
the model. The F1-score is the harmonic mean of precision and recall.
PR-AUC and ROC-AUC assess ranking performance, while the Brier score
measures the mean squared error of predicted probabilities.

\section{Experimental Results}

Table~\ref{tab:results} presents the saved-model evaluation results
on the final 15\% of the processed dataset. The reported metrics use
the thresholds stored with the respective model bundles.

\begin{table*}[t]
\centering
\caption{Comparison of Model Performance}
\label{tab:results}
\begin{tabular}{lcccccccc}
\toprule
Model & Threshold & PR-AUC & ROC-AUC & Precision & Recall &
F1-score & Accuracy & Brier Score \\
\midrule
Logistic Regression & 0.7223 & 0.3084 & 0.5682 &
0.0000 & 0.0000 & 0.0000 & 0.7331 & 0.2092 \\
XGBoost & 0.3000 & 0.3338 & 0.5985 &
0.3364 & 0.5054 & 0.4040 & 0.6019 & 0.2091 \\
\bottomrule
\end{tabular}
\end{table*}

XGBoost obtains a PR-AUC of 0.3338 and ROC-AUC of 0.5985,
compared with 0.3084 and 0.5682 for Logistic Regression.
At its evaluated threshold of 0.3000, XGBoost achieves a recall
of 0.5054 and an F1-score of 0.4040.

Logistic Regression produces no positive predictions at its saved
threshold of approximately 0.7223, resulting in zero precision,
recall, and F1-score under the evaluation convention used. Its
accuracy is 0.7331, but accuracy alone is not sufficient to assess
a model on an imbalanced classification problem.

The Brier scores are similar for the two models. These results
suggest that XGBoost performs better on the reported ranking metrics
and identifies more delayed flights at the evaluated threshold.
However, the overall discrimination performance remains limited.

\section{Discussion}

The results demonstrate the importance of selecting appropriate
evaluation metrics and classification thresholds. A high accuracy
score does not necessarily indicate effective delay detection when
the target classes are imbalanced.

XGBoost provides a more useful balance between precision and recall
at the threshold used in this evaluation. Nevertheless, the results
indicate that additional work is required before the system can be
used for operational decision-making.

Threshold selection should reflect the intended application. For
example, a system designed to identify as many delayed flights as
possible may prioritize recall, while a system intended to minimize
false alarms may prioritize precision.

\section{Limitations and Future Work}

This study has several limitations. The dataset represents flights
from a specific region and year, so the findings may not generalize
to other airports, airlines, or periods. The models also have limited
ranking performance, indicating that the current features and
approaches do not fully explain flight delays.

Future work may include improved threshold selection, further
hyperparameter tuning, richer operational and weather information,
and evaluation across different time periods. Later project phases
may investigate temporal deep learning methods and offline
rescheduling recommendations. These extensions are outside the
scope of the current Phase 1 implementation.

\section{Conclusion}

This project implements a flight-delay classification pipeline using
historical flight records, temporal feature engineering, Logistic
Regression, and XGBoost. The processed dataset contains 327,346
labeled flights, with approximately 24.47\% classified as delayed.

In the reported evaluation, XGBoost outperforms Logistic Regression
on PR-AUC and ROC-AUC and achieves a recall of approximately 0.505
at its evaluated threshold. The findings also show why accuracy
should be considered alongside precision, recall, F1-score, and
probability-based metrics. Further improvements are needed to
develop a more reliable flight-delay prediction system.

\begin{thebibliography}{3}

\bibitem{cox1958}
D. R. Cox, ``The regression analysis of binary sequences,''
\textit{Journal of the Royal Statistical Society: Series B},
vol. 20, no. 2, pp. 215--242, 1958.

\bibitem{chen2016}
T. Chen and C. Guestrin, ``XGBoost: A scalable tree boosting system,''
in \textit{Proceedings of the 22nd ACM SIGKDD International
Conference on Knowledge Discovery and Data Mining}, 2016,
pp. 785--794.

\bibitem{brier1950}
G. W. Brier, ``Verification of forecasts expressed in terms of
probability,'' \textit{Monthly Weather Review}, vol. 78, no. 1,
pp. 1--3, 1950.

\end{thebibliography}

\end{document}